\documentclass[12pt]{article}
\usepackage[margin=1in]{geometry}
\usepackage[T1]{fontenc}
\usepackage{bold-extra}
\usepackage{xcolor}
\definecolor{garnet}{HTML}{73000A}

\setlength{\parindent}{0pt}
\setlength{\parskip}{0.5em}

\begin{document}

{\color{garnet}\large\textbf{\textsc{SC Space Grant Personal Essay}}}\\[-0.8em]
{\color{garnet}\rule{\textwidth}{0.5pt}}\\[0.2em]
Growing up I found an early interest in construction. I would often find myself outside building forts or tinkering with objects in the garage. As I grew, I transitioned the outdoor construction to indoor construction with Legos. This simple hobby quickly turned into an obsession, and I would spend countless hours constructing things out of the plastic bricks.

In middle school I was given my first introduction to the world of STEM through Lego robotics. Showing me that building was not just creative but also analytical, requiring testing, refinement, and an understanding of how things work. I was able to learn basic block coding, adapt and change designs based on the challenge at hand, and develop strategic thinking based on the mission at hand. I was also introduced to the world of research, having to research technology, propose solutions to issues, and present it at competitions.

These early experiences laid the foundation for my interest in engineering and motivated me to pursue a degree in Mechanical Engineering at the University of South Carolina. As a sophomore I was drawn to undergraduate research because it allowed me to learn beyond a traditional classroom setting. Through research, I can work hands on with physical systems, develop and apply coding skills, and engage in problem solving that impacts people's lives. This experience has improved my ability to learn, my critical thinking skills, and curiosity.

In parallel, my experiences in ROTC have reinforced the skills I've developed through research while cultivating new ones. Through ROTC, I have learned to manage my time effectively, make decisions under pressure, and lead within team-based settings. Abilities that enhance both my academic and research pursuits.

Building on my current research and leadership experiences I plan to continue developing my technical and analytical skills through radar research, focusing on data processing and tracking. Upon graduation I plan to serve as an Army engineer on active duty, applying my technical expertise and leadership skills in mission driven settings. In the long term, I hope to pursue graduate school while serving as a captain. Ultimately advancing within the Army Corps of Engineers to lead engineering projects throughout the US.

These experiences will allow me to combine my passion for engineering research with leadership and service, preparing me to contribute meaningfully to both my field and the communities I serve.

\end{document}
